% Audited C-owned proof fragment, integrated in subject11 on 2026-11-19.
% Historical scope: this defines a precise finite-type category; it does not
% identify the 1964 expression `schemas morceles'.
\paragraph{A precise finite-type category of piecewise morphisms.}

Let $k$ be an algebraically closed field.  A constructible $k$-set means a
constructible subset of the $k$-points of a separated reduced finite-type
$k$-scheme.  A map is \emph{piecewise regular} if its source admits a finite
partition into locally closed reduced $k$-subschemes on each of which the map
is induced by a $k$-morphism.  Let $\mathcal P_k$ be this category.  Finite
products and their usual universal properties exist in $\mathcal P_k$.
Every object can, by using a finite affine stratification and disjoint tags,
be encoded as a constructible subset of an affine space.  The choice of
encoding does not change the category up to equivalence.

Write $\mathcal C_k$ for the category with the same objects but all maps
whose graphs are constructible.  In positive characteristic these categories
need not have the same morphisms.  Neither category records nilpotents.
An algebraic group in this subsection is a separated reduced finite-type
group scheme over $k$, and is not assumed affine or connected.

\begin{lemma}[Proved: piecewise morphisms and graphs]\label{C:piecewise-graph-category}
A piecewise regular map has constructible graph.  Conversely, in
characteristic zero every map with constructible graph is piecewise regular.
In arbitrary characteristic, a constructible graph whose projection to the
source is bijective on all scheme points and induces isomorphisms on their
residue fields is piecewise isomorphic over the source, and hence defines a
piecewise regular map.  Thus the last graph convention, interpreted on
finite-type reduced schemes and finite stratifications, gives
$\mathcal P_k$.
\end{lemma}

\begin{proof}
The graph of a morphism to a separated scheme is closed in the product of
its source with the target; a finite union over locally closed source
strata is constructible.  Composition is piecewise regular: refine the
source strata by the inverse images of strata for the second map, then
refine these constructible sets into locally closed strata.

For the converse in characteristic zero, use noetherian induction on the
closure of the source.  It suffices to find a dense open regularity locus
on each irreducible component.  Restrict to an irreducible source $V$ and
to affine source and target charts meeting the generic graph.  The graph
has one point over each point of its source.  The same uniqueness holds
after algebraically closed field extension: a failure would define a
nonempty constructible incidence set over $k$, which would already have a
$k$-point.  Consequently the dominant irreducible component of the graph
closure is generically finite over $V$ with purely inseparable function
field extension.  There is only one such component, since two dominant
components would give two points in a geometric generic fibre.  In
characteristic zero its function field is $k(V)$.  A birational morphism
of varieties is an isomorphism on suitable dense opens.  Shrinking further
puts that open of the graph closure inside the actual graph and excludes
the images of its remaining strata.  Its inverse then supplies the
required morphism on an open of $V$.  The complement has smaller dimension,
so induction gives a finite stratification.

Under the stated residue-field condition the same argument works in any
characteristic: the generic extension is already trivial.  Conversely, a
finite graph stratification by sections identifies each graph stratum
with its source stratum, so the projection is bijective on all scheme
points with identical residue fields.  The underlying facts about
radicial maps and birational open isomorphisms are precisely Stacks,
Lemmas 29.10.2 and 29.51.5 (tags 01S4 and 0BAC)\cite[tags 01S4 and 0BAC]{BStacks}.
\end{proof}

\begin{lemma}[Proved: regularity of piecewise homomorphisms]\label{C:piecewise-group-fullness}
Every piecewise regular homomorphism $f:G(k)\to H(k)$ between algebraic
$k$-groups is induced by a unique homomorphism of algebraic groups.
\end{lemma}

\begin{proof}
Remove the other irreducible components from one irreducible component
of $G$, and then choose a nonempty open regularity stratum for $f$.
This gives a $V$ open in all of $G$ on which $f$ is regular.
Choose $v\in V(k)$.  For each $g\in G(k)$ the open
$g v^{-1}V$ contains $g$.  On this open the homomorphism identity gives
\[
 f(z)=f(gv^{-1})\,f(vg^{-1}z).
\]
The right-hand side is a morphism.  These morphisms agree on overlaps,
since they agree on all $k$-points and the schemes are reduced.  They glue
to a morphism on $G$.  The group identities hold as scheme morphisms by
the same reducedness argument, since products of reduced finite-type
schemes over the algebraically closed field $k$ are reduced.
Uniqueness also follows from density of
closed points.  No restriction on the characteristic is used here.
\end{proof}

\begin{theorem}[Proved: algebraization in the specified finite-type category]\label{C:piecewise-group-algebraization}
The functor from algebraic $k$-groups to group objects in $\mathcal P_k$ is
an equivalence of categories.  In particular, if $\operatorname{char}k=0$,
it is an equivalence with group objects in $\mathcal C_k$.
\end{theorem}

\begin{proof}
Full faithfulness is the preceding lemma.  We prove essential surjectivity.
Let $D$ be a constructible set with group operations that are piecewise
regular.  Encode $D$ in an affine space and put $d=\dim D$.  If $d=0$,
then $D$ is finite and the constant finite group scheme proves the claim.
Assume $d>0$.

\emph{Step 1: a normal finite-index subgroup with one generic component.}
Choose disjoint constructible subsets $C_1,\ldots,C_m$ of $D$, each dense
open in one of the distinct $d$-dimensional irreducible components of
$\overline D$, such that the complement of their union has dimension less
than $d$.  Two constructible subsets are identified here if their
symmetric difference has dimension less than $d$.  Modulo this relation,
the $C_i$ are the atoms of the finite Boolean algebra of constructible
subsets of $D$.

A piecewise regular bijection with piecewise regular inverse preserves
dimension and these generic atoms.  Hence left translation gives an
action of $D$ on $\{C_1,\ldots,C_m\}$.  The action is constructible: for
fixed $i,j$ its relevant fibre is
\[
 \{g\in D:\dim(gC_i\cap C_j)=d\}.
\]
Constructibility follows by stratifying the incidence correspondence and
applying constructibility of fibre dimension (Stacks, Lemma 37.30.3,
tag 05F9)\cite[tag 05F9]{BStacks}.  This applies to piecewise regular maps by a finite
stratification into ordinary morphisms.

The action is transitive.  Indeed, the piecewise regular map
\[
 C_j\times C_i\longrightarrow D,\qquad (b,a)\longmapsto ba^{-1}
\]
has a source of dimension $2d$.  Its fibre over $g$ is isomorphic as a
constructible set to $\{a\in C_i:ga\in C_j\}$.  If every fibre had
dimension at most $d-1$, the dimension-of-fibres inequality would bound
the source dimension by $2d-1$, a contradiction.  A fibre of dimension
$d$ shows that some $g$ carries the generic atom of $C_i$ to that of $C_j$.

Let $D^0$ be the stabilizer of $C_1$.  It is constructible, has index $m$,
and has dimension $d$.  If it had $s$ generic atoms, each of its $m$
cosets would have $s$ atoms, and the disjoint coset partition would give
$m=ms$.  Thus $s=1$.  A constructible subgroup of finite index in $D^0$
must therefore be all of $D^0$.  Applying this to
$D^0\cap gD^0g^{-1}$ shows that $D^0$ is normal in $D$.

\emph{Step 2: obtain a strict birational group law.}
Let $V$ be the unique $d$-dimensional generic component of $D^0$.
Choose a dense open $V_0\subset V$ contained in $D^0$.  Every translation
of $D^0$ preserves its unique generic atom.  In particular the product
of two independent generic points lies in that atom: for each fixed
$a$, the bad set of $b$ for which $ab$ lies outside $V_0$ has dimension
less than $d$, by translation invariance; the fibre-dimension inequality
therefore bounds the total bad set in $D^0\times D^0$ by $2d-1$.
The multiplication,
left-division and right-division maps are piecewise regular, so the three
pairwise projections of the generic multiplication graph induce
birational maps to $V\times V$.  More explicitly,
$(a,b)\mapsto(a,ab)$ and $(a,b)\mapsto(ab,b)$ have the piecewise regular
inverses $(a,c)\mapsto(a,a^{-1}c)$ and
$(c,b)\mapsto(cb^{-1},b)$.  On dense opens these are inverse rational
maps, and all three coordinates lie in $V_0$.

Take a dense open $W_0$ of the closure of this graph which is contained
in the actual multiplication graph, lies in $V_0^3$, and on which all
three pairwise projections are open immersions.  Denote their dense open
images in $V^2$ by $U_{12},U_{13},U_{23}$.  The projections of each
$U_{ij}$ onto either factor are dense opens of $V$.  Let $X\subset V_0$
be their intersection, and set $W=W_0\cap X^3$.

We check the fibre-density condition, which cannot be omitted.  Fix any
one coordinate $a\in X$.  The fibre of $W_0$ at this coordinate maps
isomorphically, by either remaining coordinate, onto a nonempty open of
$V$: the two relevant pairwise projections are open immersions, and
$a$ belongs to both projection images by the choice of $X$.  Thus it
is an irreducible open variety, birational to $V$.  Requiring the other
two coordinates to lie in $X$ removes only proper closed subsets of
this fibre.  The resulting fibre is nonempty, and both of its coordinate
images are dense opens of $X$.  It follows that every pairwise projection
of $W$ is an open immersion into $X^2$ with dense fibres over either
factor.  Associativity on its domain is inherited from $D^0$.
Consequently $(X,W)$ satisfies exactly Definition 2.1 of
Edixhoven--Romagny\cite[Definition 2.1, pp.~4--5]{CEdixhovenRomagny2013}, not merely a weaker generic associativity condition.

\emph{Step 3: algebraize and identify the entire connected part.}
Apply Edixhoven--Romagny, Theorem 2.11, whose complete proof is in their
Section 2\cite[Theorem 2.11 and its proof, pp.~4--7]{CEdixhovenRomagny2013}.  It constructs an algebraic group $A$ containing $X$ as a
dense open.  Its underlying abstract group is the subgroup of
$\operatorname{Bir}(X)$ generated by the rational left translations
$\lambda_x$ for $x\in X(k)$.  It is separated and of finite type, since
$X$ is of finite type.  It is irreducible because $X$ is irreducible and
dense.

For every $h\in D^0$, left translation by $h$ gives a birational
self-map of $V$, hence an element $\rho(h)$ of
$\operatorname{Bir}(X)$.  The map $\rho$ is a homomorphism and is
injective: if $ha=a$ on a nonempty open, cancellation gives $h=1$.
Moreover $D^0=X X^{-1}$.  To see this, for any $h$ choose
$y\in X\cap h^{-1}X$, which is nonempty since both sets have the unique
generic atom, and write $h=(hy)y^{-1}$.  Thus $\rho(D^0)$ is precisely
the subgroup generated by the $\lambda_x$, namely $A(k)$.

This isomorphism $\rho:D^0\to A(k)$ is constructible: its graph is the
image of the constructible map
\[
 X\times X\longrightarrow D^0\times A,
 \qquad (x,y)\longmapsto
 \bigl(xy^{-1},\lambda_x\lambda_y^{-1}\bigr).
\]
We must also check piecewise regularity in arbitrary characteristic.
The translates of the dense open $X\subset A$ cover $A$, and
quasi-compactness of the finite-type scheme $A$ supplies a finite
subcover $\rho(h_i)X$.  Pulling this cover back through the already
established abstract group isomorphism gives $D^0=\bigcup_i h_iX$.
On $h_iX$, the map $\rho$ is the composite of the
piecewise regular translation $h_iX\to X$ with the regular map
$x\mapsto\rho(h_i)\lambda_x$.  Hence $\rho$ is piecewise regular.
The corresponding opens $\rho(h_i)X$ cover $A$, and on each of them the
inverse is a regular chart followed by a piecewise regular translation
in $D^0$.  Therefore $\rho$ is an isomorphism in $\mathcal P_k$.

\emph{Step 4: restore the finitely many components.}
Choose left coset representatives $r_i$ for $D/D^0$.  On each coset
$r_iD^0$ transport the variety structure of $A$ by left translation.
For $r_ir_jD^0=r_\ell D^0$ put
$c_{ij}=r_\ell^{-1}r_ir_j\in D^0$.  In the transported coordinates,
multiplication is
\[
 (r_i a)(r_j b)=r_\ell c_{ij}(r_j^{-1}ar_j)b.
\]
Conjugation by $r_j$ is a piecewise regular automorphism of $D^0$;
after transport to $A$ it is regular by
Lemma~\ref{C:piecewise-group-fullness}.  Thus multiplication is regular
on every pair of components.  The same formula applied to inverses
shows that inversion is regular.  The finite disjoint union is a
separated reduced finite-type algebraic group, and its identification
with $D$ is piecewise regular in both directions.  This proves essential
surjectivity and the equivalence.  Characteristic zero follows from
Lemma~\ref{C:piecewise-graph-category}.
\end{proof}

\paragraph{Positive-characteristic distinction.}
For $\operatorname{char}k=p>0$, the inverse of Frobenius on
$\mathbb G_a(k)$ is an additive homomorphism with closed graph
$\{(x,y):y^p=x\}$.  It is not regular: a polynomial $P(x)$ satisfying
$P(x)^p=x$ cannot exist by degree.  It is not even rational on a dense
open, since $x$ is not a $p$th power in $k(x)$ (use its valuation at
$x=0$).  It therefore does not belong to $\mathcal P_k$.  Its graph
projection has generic residue-field extension $k(x)\subset k(y)$,
$x=y^p$, of degree $p$, so it also fails the residue-field-trivial
convention.  The ordinary functor from algebraic groups to group objects
in $\mathcal C_k$ is consequently not full in characteristic $p$.
The theorem above makes no unproved assertion that all objects of that
larger characteristic-$p$ category are algebraizable; the point here
is the proved equivalence for $\mathcal P_k$ and the precise obstruction
to replacing its morphisms by arbitrary constructible graphs.

\paragraph{Historical limitation.}
Hrushovski's \emph{Valued fields, metastable groups}, introduction,
pp.~1--2\cite[pp.~1--2 and p.~2, footnote 3]{CHrushovskiMetastableDraft}, records both the characteristic-zero equivalence and the
residue-field-trivial convention, attributing the latter to Grothendieck's
letter to Poizat of 17 June 1979.  Poizat's \emph{Logique et g\'eom\'etrie},
Sections 2.B and 3.B\cite[Sections 2.B and 3.B]{CPoizat2025}, spells out the characteristic distinction.  These
opened sources justify the specified category.  They do not identify
it with every scheme, pro-affine object, sheaf or partition construction
that might have been intended by the 1964 expression ``sch\'emas
morcel\'es''.  That historical identification remains a separate task.
